\ifdefined\gkver\else\def\gkver{student}\fi
\documentclass[\gkver]{gaokaozhenti}
\begin{document}

\chapter{1990年上海}

\begin{enumerate}
\item
%% number: 1
%% paperName: 1990年上海高考第 1 题 4 分
%% typeId: 1
%% type: 单选题
%% chapter: 光学
%% point: 光的电磁说
%% method: 无
%% score: 4
%% degree: 850
%% duplicateId: 0
%% body:
单色光从空气射到水中，它的\xzanswer{B}
\fourchoices[answer=B]
{频率和波长都要改变}
{波长和传播速度都要改变}
{传播速度和颜色都要改变}
{频率和颜色都要改变}
%% answer: B
%% memo:
\memoanswer{
光在传播的过程中，频率不变，频率不变即颜色不变，但波长和传播速度都要改变，故 B 正确。
}

\item
%% number: 2
%% paperName: 1990年上海高考第 2 题 4 分
%% typeId: 1
%% type: 单选题
%% chapter: 固体液体的性质
%% point: 固体的基本性质
%% method: 无
%% score: 4
%% degree: 850
%% duplicateId: 0
%% body:
把薄片的一面涂上一薄层石蜡，然后用烧热的钢针接触它的反面，熔化了的石蜡呈椭圆形，那么，这薄片是\xzanswer{C}
\fourchoices[answer=C]
{非晶体}
{多晶体}
{单晶体}
{无法判定}
%% answer: C
%% memo:
\memoanswer{
由于熔化的石蜡呈椭圆形，即各个方向的导热性能不同，表现为各向异性，所以这薄片为单晶体，故 C 正确。
}

\item
%% number: 3
%% paperName: 1990年上海高考第 3 题 4 分
%% typeId: 1
%% type: 单选题
%% chapter: 功和能
%% point: 动能定理
%% method: 无
%% score: 4
%% degree: 650
%% duplicateId: 0
%% body:
汽车在平直公路上行驶，在它的速度从零增加到 $v$ 的过程中，汽车发动机做的功为 $W_{1}$；在它的速度从 $v$ 增加到 $2v$ 的过程中，汽车发动机做的功为 $W_{2}$；设汽车在行驶过程中发动机的牵引力和所受阻力都不变，则有\xzanswer{B}
\fourchoices[answer=B]
{$W_{2}=2W_{1}$}
{$W_{2}=3W_{1}$}
{$W_{2}=4W_{1}$}
{仅能判定 $W_{2}>W_{1}$}
%% answer: B
%% memo:
\memoanswer{
由题意知，汽车牵引力 $F$ 和所受阻力不变，由做功的定义知 $W_{1}=Fs_{1}$，$W_{2}=Fs_{2}$。

由动能定理得 $(F-f)s_{1}=\frac{1}{2}mv^{2}-0$，$(F-f)s_{2}=\frac{1}{2}m(2v)^{2}-\frac{1}{2}mv^{2}$，得 $s_{2}=3s_{1}$，所以 $W_{2}=3W_{1}$，故 B 正确。
}

\item
%% number: 4
%% paperName: 1990年上海高考第 4 题 4 分
%% typeId: 1
%% type: 单选题
%% chapter: 原子和原子核
%% point: 半衰期
%% method: 无
%% score: 4
%% degree: 650
%% duplicateId: 0
%% body:
设某放射性同位素 A 的半衰期为 $T$，另一种放射性同位素 B 的半衰期为 $\frac{T}{2}$。在初始时刻，A 的原子核数目为 $N_{0}$，B 的原子核数目为 $4N_{0}$，则\xzanswer{B}
\fourchoices[answer=B]
{经过时间 $T$，A、B 的原子核数目都等于 $\frac{N_{0}}{2}$}
{经过时间 $2T$，A、B 的原子核数目都等于 $\frac{N_{0}}{4}$}
{经过时间 $3T$，A、B 的原子核数目都等于 $\frac{N_{0}}{8}$}
{经过时间 $4T$，A、B 的原子核数目都等于 $\frac{N_{0}}{16}$}
%% answer: B
%% memo:
\memoanswer{
由半衰期定义有 $N_{A}=N_{0}\cdot\left(\frac{1}{2}\right)^{\frac{t}{T}}$，$N_{B}=4N_{0}\cdot\left(\frac{1}{2}\right)^{\frac{2t}{T}}$。

逐项计算结果，A 项：$N_{A}=\frac{N_{0}}{2}$，$N_{B}=N_{0}$；B 项：$N_{A}=\frac{N_{0}}{4}$，$N_{B}=\frac{N_{0}}{4}$；C 项：$N_{A}=\frac{N_{0}}{8}$，$N_{B}=\frac{N_{0}}{16}$；D 项：$N_{A}=\frac{N_{0}}{16}$，$N_{B}=\frac{N_{0}}{64}$，故 B 正确。
}

\item
%% number: 5
%% paperName: 1990年上海高考第 5 题 4 分
%% typeId: 1
%% type: 单选题
%% chapter: 气体性质
%% point: 查理定律
%% method: 无
%% score: 4
%% degree: 650
%% duplicateId: 0
%% body:
如图所示密封的 U 形管中装有水银，左右两端都封有空气，两水银面高度差为 $h$。把 U 形管竖直浸没在热水中，高度差 $h$ 将\xzanswer{A}
\onepicture[w=2cm]{\includesvg{figs/05}}
\fourchoices[answer=A]
{增大}
{减小}
{不变}
{两侧空气柱的长度未知，不能判断}
%% answer: A
%% memo:
\memoanswer{
设温度为 $T_{0}$ 时，左侧空气的压强为 $p_{1}$，右侧空气的压强为 $p_{2}$，由压强平衡得 $p_{2}=p_{1}+h$，所以 $p_{2}>p_{1}$。现温度升高 $\Delta T$，假设两侧空气体积不变，即气体作等容变化，左侧空气的压强为 $p_{1}'$，右侧空气的压强为 $p_{2}'$，则有 $\frac{p_{1}}{T}=\frac{p_{1}'}{T+\Delta T}$，解得 $p_{1}'=(1+\frac{\Delta T}{T})p_{1}$，则 $\Delta p_{1}=\frac{\Delta T}{T}p_{1}$，同理有 $\Delta p_{2}=\frac{\Delta T}{T}p_{2}$，故 $\Delta p_{2}>\Delta p_{1}$，则右侧水银液面下降，水银的高度差将增大，故 A 正确。

\onepicture[w=4.5cm]{figs/05_an.jpg}
}

\item
%% number: 6
%% paperName: 1990年上海高考第 6 题 4 分
%% typeId: 1
%% type: 单选题
%% chapter: 力物体的平衡
%% point: 平衡综合
%% method: 无
%% score: 4
%% degree: 650
%% duplicateId: 0
%% body:
如图所示，质量为 $m$ 的匀质木杆，上端可绕固定水平光滑轴 O 转动，下端搁在木板上，木板置于光滑水平地面，棒与竖直线成 $45^{\circ}$ 角，棒与木板间的摩擦系数为 0.5。为使木板向右作匀速运动，水平拉力 $F$ 等于\xzanswer{D}
\onepicture[align=right, w=5cm]{\includesvg{figs/06}}
\fourchoices[answer=D]
{$\frac{1}{2}mg$}
{$\frac{1}{3}mg$}
{$\frac{1}{4}mg$}
{$\frac{1}{6}mg$}
%% answer: D
%% memo:
\memoanswer{
对木杆，木板对杆的支持力为 $N$，摩擦力为 $f$，由力矩平衡得 $NL\sin45^{\circ}+fL\cos45^{\circ}=mg\frac{L}{2}\sin45^{\circ}$，又 $f=\mu N=\frac{1}{2}N$，联立解得 $N=\frac{1}{3}mg$，$f=\frac{1}{6}mg$，由木板的平衡条件得 $F=f=\frac{1}{6}mg$，故 D 正确。
}

\item
%% number: 7
%% paperName: 1990年上海高考第 7 题 4 分
%% typeId: 1
%% type: 单选题
%% chapter: 电磁感应
%% point: 楞次定律推广
%% method: 无
%% score: 4
%% degree: 450
%% duplicateId: 0
%% body:
如图所示，ab 是一个可绕垂直于纸面的轴 O 转动的闭合框，当滑线变阻器 $R$ 的滑片 P 自左向右滑行时，线框 ab 将\xzanswer{C}
\onepicture[w=5.5cm]{\includesvg{figs/07}}
\fourchoices[answer=C]
{保持静止不动}
{逆时针转动}
{顺时针转动}
{发生转动，但因电源极性不明，无法确定转动方向}
%% answer: C
%% memo:
\memoanswer{
设电源左正右负，因 $R$ 变大，$I$ 变小，磁通量 $\Phi$ 减小，感应电流 a 端向里、b 端向外，在安培力力矩作用下，线框顺时针转动；设电源左负右正，也是线框顺时针转动，故 C 正确。
}

\item
%% number: 8
%% paperName: 1990年上海高考第 8 题 4 分
%% typeId: 1
%% type: 单选题
%% chapter: 磁场
%% point: 磁感应强度
%% method: 无
%% score: 4
%% degree: 650
%% duplicateId: 0
%% body:
磁感应强度的单位是特斯拉，$1\UT$ 相当于\xzanswer{A}
\fourchoices[answer=A]
{$1\Ukg/(\UA\cdot\Us^{2})$}
{$1\Ukg\cdot\Um/(\UA\cdot\Us^{3})$}
{$1\Ukg\cdot\Um^{2}/\Us^{2}$}
{$1\Ukg\cdot\Um/(\UA\cdot\Us^{2})$}
%% answer: A
%% memo:
\memoanswer{
由磁感应强度定义式知 $B=\frac{F}{IL}$，故 $1\UT=1\frac{\UN}{\UA\cdot\Um}$，又 $1\UT=1\frac{\UN}{\UA\cdot\Um}=1\frac{\Ukg\cdot\Um/\Us^{2}}{\UA\cdot\Um}=1\frac{\Ukg}{\UA\cdot\Us^{2}}$，故 A 正确。
}

\item
%% number: 9
%% paperName: 1990年上海高考第 9 题 5 分
%% typeId: 2
%% type: 多选题
%% chapter: 原子和原子核
%% point: 玻尔模型
%% method: 无
%% score: 5
%% degree: 650
%% duplicateId: 0
%% body:
按照玻尔理论，在氢原子中，当电子从半径为 $4r_{1}$ 的轨道跃迁到半径为 $r_{1}$ 的轨道时，它的能量变化是\xzanswer{AD}
\fourchoices[answer=AD]
{电势能减少，动能增加}
{电势能减少，动能减少}
{电势能的减少等于动能的增加}
{电势能的减少大于动能的增加}
%% answer: AD
%% memo:
\memoanswer{
因 $E_{n}=\frac{-13.6\UeV}{n^{2}r_{1}^{2}}$，$n$ 从 4 到 1，总能量减少。因库仑力提供向心力有 $k\frac{e^{2}}{r^{2}}=m\frac{v^{2}}{r}$，则 $E_{k}=\frac{ke^{2}}{2r}$，而 $E_{p}=-\frac{ke^{2}}{r}$，所以 $E_{p}=-2E_{k}$，又 $E_{p}+E_{k}=E$，所以电势能 $E_{p}=\frac{-27.2\UeV}{n^{2}r_{1}^{2}}$，动能 $E_{k}=\frac{13.6\UeV}{n^{2}r_{1}^{2}}$，故 $n$ 从 4 到 1，电势能减少、动能增加，并且电势能的减少大于动能的增加，故 AD 正确。
}

\item
%% number: 10
%% paperName: 1990年上海高考第 10 题 5 分
%% typeId: 2
%% type: 多选题
%% chapter: 交流电
%% point: LC振荡回路
%% method: 无
%% score: 5
%% degree: 650
%% duplicateId: 0
%% body:
如图所示的 LC 振荡电路，当电键 K 打向右边发生振荡后，下列说法中正确的是\xzanswer{ABD}
\onepicture[w=4.5cm]{\includesvg{figs/10}}
\fourchoices[answer=ABD]
{振荡电流达到最大值时，电容器上的带电量为零}
{振荡电流达到最大值时，磁场能最大}
{振荡电流为零时，电场能为零}
{振荡电流相邻两次为零的时间间隔等于振荡周期的一半}
%% answer: ABD
%% memo:
\memoanswer{
由振荡电路的特性知，当振荡电流达到最大值时，此时的磁场能最大，故电场能最小，所以电容器上的带电量为零，故 A、B 正确；当振荡电流为零时，此时的磁场能为零，故电场能最大，所以电容器上的带电量达到最大，故 C 错误；振荡电流相邻两次为零的时间间隔等于振荡周期的一半，故 D 正确。

\onepicture[w=6cm]{figs/10_an.jpg}
}

\item
%% number: 11
%% paperName: 1990年上海高考第 11 题 5 分
%% typeId: 2
%% type: 多选题
%% chapter: 机械波
%% point: 波的图象
%% method: 无
%% score: 5
%% degree: 650
%% duplicateId: 0
%% body:
如图所示为某一时刻简谐横波的图像，波的传播方向沿 $x$ 正方向，下列说法正确的是\xzanswer{AD}
\onepicture[w=5.5cm]{\includesvg{figs/11}}
\fourchoices[answer=AD]
{质点 A、D 的振幅相等}
{在该时刻质点 B、E 的速度大小和方向相同}
{在该时刻质点 C、F 的加速度为零}
{在该时刻质点 D 正向下运动}
%% answer: AD
%% memo:
\memoanswer{
由波动图象可知，所有质点振幅应该相同，所以 A 正确。

根据微平移法，如图所示下一时刻的波形，所以 B 点向上振动，E 点向下振动，D 点向下振动，所以 B 错误、D 正确。

\onepicture[w=5.5cm]{\includesvg{figs/11_an}}

由于此刻 C、F 的位移最大，所以回复力最大，因此加速度最大，C 错误。
}

\item
%% number: 12
%% paperName: 1990年上海高考第 12 题 5 分
%% typeId: 2
%% type: 多选题
%% chapter: 运动和力
%% point: 动量守恒定律
%% method: 无
%% score: 5
%% degree: 650
%% duplicateId: 0
%% body:
A、B 两球在光滑水平面上沿同一直线、同一方向运动，A 球的动量是 $5\Ukg\cdot\Ums$，B 球的动量是 $7\Ukg\cdot\Ums$，当 A 球追上 B 球时发生碰撞，则碰撞后 A、B 两球的动量可能值是\xzanswer{BC}
\fourchoices[answer=BC]
{$p_{A}=6\Ukg\cdot\Ums$，$p_{B}=6\Ukg\cdot\Ums$}
{$p_{A}=3\Ukg\cdot\Ums$，$p_{B}=9\Ukg\cdot\Ums$}
{$p_{A}=-2\Ukg\cdot\Ums$，$p_{B}=14\Ukg\cdot\Ums$}
{$p_{A}=-6\Ukg\cdot\Ums$，$p_{B}=15\Ukg\cdot\Ums$}
%% answer: BC
%% memo:
\memoanswer{
两球碰撞，内力远大于外力，满足动量守恒，由动量守恒可以判断前三个选项都符合，故 D 错误；由已知条件得 $p_{A}<p_{B}$，$v_{A}>v_{B}$（A 追上 B），故 $m_{A}<m_{B}$。A 项中因 $p_{A}'=p_{B}'$，所以 $v_{A}'>v_{B}'$，A 要越过 B，是不可能的，同理判断 B、C、D 项的情况可能存在，故 A 错误；又因碰撞时可能存在能量损失，故系统的总能量不会增加，由动量与动能的关系有 $E_{k}=\frac{p^{2}}{2m}$，

系统的初动能为 $E_{\text{初}}=\frac{25}{2m_{A}}+\frac{49}{2m_{B}}$。B 项中系统的末动能为 $E_{B}=\frac{9}{2m_{A}}+\frac{81}{2m_{B}}$，则 $E_{\text{初}}-E_{B}=\frac{16}{2m_{A}}-\frac{32}{2m_{B}}$ 可能小于等于 0，故 B 正确；C 项中系统的末动能为 $E_{C}=\frac{4}{2m_{A}}+\frac{196}{2m_{B}}$，则 $E_{\text{初}}-E_{C}=\frac{5}{2m_{A}}-\frac{115}{2m_{B}}$ 可能小于等于 0，故 C 正确；D 项中系统的末动能为 $E_{D}=\frac{36}{2m_{A}}+\frac{225}{2m_{B}}$，则 $E_{\text{初}}-E_{D}=\frac{-27}{2m_{A}}-\frac{144}{2m_{B}}<0$，即系统的总能量增加，故 D 错误。
}

\item
%% number: 13
%% paperName: 1990年上海高考第 13 题 5 分
%% typeId: 2
%% type: 多选题
%% chapter: 电路
%% point: 动态分析
%% method: 等效替代
%% score: 5
%% degree: 350
%% duplicateId: 0
%% body:
如图所示电路，总电压 $U$ 保持不变，滑线变阻器总电阻为 $2R$。当滑片位于变阻器中点 O 时，四个电流表 A$_{1}$、A$_{2}$、A$_{3}$、A$_{4}$ 上读数相等，都等于 $I_{0}$。当滑片向上移动到 O$'$ 点时\xzanswer{BC}
\onepicture[w=3.5cm]{\includesvg{figs/13}}
\fourchoices[answer=BC]
{A$_{1}$ 的示数大于 $I_{0}$}
{A$_{2}$ 的示数大于 $I_{0}$}
{A$_{3}$ 的示数大于 $I_{0}$}
{A$_{4}$ 的示数大于 $I_{0}$}
%% answer: BC
%% memo:
\memoanswer{
滑片位于中点 O 时，电桥平衡，4 个电流表示数相等。当滑片移到 O$'$ 时，A$_{3}$ 支路电阻变小，A$_{4}$ 支路电阻变大，所以 A$_{1}$ 和 A$_{3}$ 所在支路电阻并联值小于 A$_{2}$ 和 A$_{4}$ 所在支路电阻并联值，因两部分串联分压，所以 A$_{1}$ 所在支路电压低于原来的电压，A$_{2}$ 所在支路电压高于原来，所以 A$_{1}$ 电流小于 $I_{0}$，A$_{2}$ 电流大于 $I_{0}$，故 A 错误、B 正确；滑片移到 O$'$ 时，总电阻变大，所以总电流变小，因此 A$_{4}$ 中电流变小，故 D 错误；初始时 $I_{0}R_{3}+I_{0}R_{4}=U$，$R_{3}=R_{4}=R$，联立有 $2I_{0}R=U$，变化后有 $I_{3}R_{3}'+I_{4}R_{4}'=U$ ①，$R_{3}'+R_{4}'=2R$，假设 $I_{4}$ 不变，为 $I_{0}$，则由以上各式联立推出 $I_{3}=I_{0}$，现 $I_{4}$ 变小，由①式推出 $I_{3}$ 变大，所以 A$_{3}$ 示数大于 $I_{0}$，故 C 正确。
}

\item
%% number: 14
%% paperName: 1990年上海高考第 14 题 4 分
%% typeId: 3
%% type: 填空题
%% chapter: 气体性质
%% point: 气体图象
%% method: 无
%% score: 4
%% degree: 650
%% duplicateId: 0
%% body:
一定质量的理想气体，从如图所示的 $p-T$ 图上的状态 A 变化到状态 B。在此过程中，气体分子的平均动能 \tkanswer{减小}；气体体积 \tkanswer{减小}。（填“增大”、“减小”或“不变”）
\onepicture[align=right, w=4cm]{\includesvg{figs/14}}
%% answer: 减小、减小
\jdanswer{
减小、减小
}
%% memo:
\memoanswer{
由 $p-T$ 图象知，$T_{B}<T_{A}$，因此气体分子的平均动能减小；在 $p-T$ 图象上，OB 的斜率大于 OA 的斜率，即 $\frac{p_{B}}{T_{B}}>\frac{p_{A}}{T_{A}}$，而 $\frac{p_{B}V_{B}}{T_{B}}=\frac{p_{A}V_{A}}{T_{A}}$，得 $V_{B}<V_{A}$。

\onepicture[w=4cm]{figs/14_an.jpg}
}

\item
%% number: 15
%% paperName: 1990年上海高考第 15 题 4 分
%% typeId: 3
%% type: 填空题
%% chapter: 磁场
%% point: 带电粒子在磁场中的运动
%% method: 无
%% score: 4
%% degree: 650
%% duplicateId: 0
%% body:
一初速为零的带电粒子经过电压为 $U$ 的电场加速后垂直进入磁感应强度为 $B$ 的匀强磁场中，已知带电粒子的质量是 $m$、电量是 $q$，则带电粒子所受的洛伦兹力为 \tkanswer[2.8cm]{$qB\sqrt{\frac{2qU}{m}}$}，轨道半径为 \tkanswer[2.8cm]{$\frac{1}{B}\sqrt{\frac{2mU}{q}}$}。
%% answer: $qB\sqrt{\frac{2qU}{m}}$；$\frac{1}{B}\sqrt{\frac{2mU}{q}}$
\jdanswer{
$qB\sqrt{\frac{2qU}{m}}$，$\frac{1}{B}\sqrt{\frac{2mU}{q}}$
}
%% memo:
\memoanswer{
带电粒子在电场中运动，由动能定理得 $qU=\frac{1}{2}mv^{2}$，解得 $v=\sqrt{\frac{2qU}{m}}$，则洛伦兹力 $f=qvB=qB\sqrt{\frac{2qU}{m}}$；由洛伦兹力提供向心力得 $f=qvB=m\frac{v^{2}}{r}$，解得 $r=\frac{mv}{qB}=\frac{1}{B}\sqrt{\frac{2mU}{q}}$。
}

\item
%% number: 16
%% paperName: 1990年上海高考第 16 题 4 分
%% typeId: 3
%% type: 填空题
%% chapter: 曲线运动
%% point: 人造地球卫星
%% method: 无
%% score: 4
%% degree: 450
%% duplicateId: 0
%% body:
已知地球的半径为 $R$，自转角速度为 $\omega$，地球表面的重力加速度为 $g$，在赤道上空一颗相对地球静止的同步卫星离开地面的高度是 \tkanswer[3cm]{$\sqrt[3]{\frac{gR^{2}}{\omega^{2}}}-R$}（用以上三个量表示）。
%% answer: $\sqrt[3]{\frac{gR^{2}}{\omega^{2}}}-R$
\jdanswer{
$\sqrt[3]{\frac{gR^{2}}{\omega^{2}}}-R$
}
%% memo:
\memoanswer{
由万有引力提供向心力和牛顿第二定律得 $G\frac{Mm}{(R+h)^{2}}=m\omega^{2}(R+h)$，又 $G\frac{Mm}{R^{2}}=mg$，联立解得 $h=\sqrt[3]{\frac{gR^{2}}{\omega^{2}}}-R$。
}

\item
%% number: 17
%% paperName: 1990年上海高考第 17 题 4 分
%% typeId: 7
%% type: 作图题
%% chapter: 光学
%% point: 透镜
%% method: 无
%% score: 4
%% degree: 650
%% duplicateId: 0
%% body:
如图所示，O$_{1}$O$_{2}$ 为透镜的主轴，S 为点光源，S$'$ 是 S 的像。试由作图法在图中画出透镜光心位置 O 和焦点位置 F。
\drawpicanswer{\onepicture[align=right, w=5.5cm]{\includesvg{figs/17}}}{\onepicture[align=right, w=6cm]{\includesvg{figs/17_an}}}
%% answer: 如图所示
\jdanswer{
作图如图所示。
}
%% memo:
\memoanswer{
利用作图法确定透镜的光心位置和焦点位置。

①连接 $S'S$，延长交 $O_{1}O_{2}$ 于 O，O 为光心；

②过 O 作 L 垂直于 $O_{1}O_{2}$，L 为透镜；

③过 S 作 $SC\parallel O_{1}O_{2}$ 交 L 于 C，延长 $S'C$ 交 $O_{1}O_{2}$ 于 $F_{1}$，$F_{1}$ 为焦点；

④在 $OO_{1}$ 上作 $F_{2}$，使 $OF_{2}=OF_{1}$，则 $F_{2}$ 为另一焦点。
}

\item
%% number: 18
%% paperName: 1990年上海高考第 18 题 4 分
%% typeId: 3
%% type: 填空题
%% chapter: 直线运动
%% point: 自由落体运动
%% method: 无
%% score: 4
%% degree: 650
%% duplicateId: 0
%% body:
一矿井深为 $125\Um$，在井口每隔一定时间自由下落一个小球。当第 11 个小球刚从井口开始下落时，第 1 个小球恰好到达井底，则相邻两个小球开始下落的时间间隔为 \tkanswer[1.5cm]{$0.50\Us$}，这时第 8 个小球和第 5 个小球相隔 \tkanswer[1.5cm]{$35\Um$}。
%% answer: 0.50 s；35 m
\jdanswer{
$0.50\Us$，$35\Um$
}
%% memo:
\memoanswer{
由运动学公式得 $h=\frac{1}{2}gt^{2}$，解得 $t=5\Us$，相邻两个小球开始下落的时间间隔 $\Delta t=\frac{t}{10}=0.5\Us$；当第 11 个小球刚从井口开始下落时，第 3 个小球已下落 $4\Us$，第 5 个小球已下落 $3\Us$，两球相距的距离为 $\Delta h=\frac{1}{2}g(4^{2}-3^{2})\Um=35\Um$。
}

\item
%% number: 19
%% paperName: 1990年上海高考第 19 题 4 分
%% typeId: 3
%% type: 填空题
%% chapter: 功和能
%% point: 动能定理
%% method: 无
%% score: 4
%% degree: 650
%% duplicateId: 0
%% body:
质量为 $m$ 的物体从高为 $h$ 的斜面顶端自静止开始滑下，最后停在平面上的 B 点，如图所示。若该物体从斜面顶端以初速 $v_{0}$ 沿斜面滑下，则停在平面上的 C 点。已知 AB $=$ BC，则物体在斜面上克服摩擦力所作的功为 \tkanswer[3cm]{$mgh-\frac{1}{2}mv_{0}^{2}$}。
\onepicture[align=right, w=6cm]{\includesvg{figs/19}}
%% answer: $mgh-\frac{1}{2}mv_{0}^{2}$
\jdanswer{
$mgh-\frac{1}{2}mv_{0}^{2}$
}
%% memo:
\memoanswer{
物体运动的过程中，由动能定理得 $mgh-W_{f}-\mu mg\cdot\overline{AB}=0$，$\frac{1}{2}mv_{0}^{2}+mgh-W_{f}-\mu mg\cdot2\overline{AB}=0$，联立解得 $W_{f}=mgh-\frac{1}{2}mv_{0}^{2}$。
}

\item
%% number: 20
%% paperName: 1990年上海高考第 20 题 4 分
%% typeId: 3
%% type: 填空题
%% chapter: 光学
%% point: 光的折射
%% method: 无
%% score: 4
%% degree: 650
%% duplicateId: 0
%% body:
有人在游泳池边上竖直向下观察池水的深度，看上去池水的视深约为 $h$。已知水的折射率 $n=\frac{4}{3}$，那么，水的实际深度约为 \tkanswer[2cm]{$\frac{4}{3}h$}。
%% answer: $\frac{4}{3}h$
\jdanswer{
$\frac{4}{3}h$
}
%% memo:
\memoanswer{
由题意作光路图如图所示。由折射定律得 $n=\frac{\sin i}{\sin r}$（$i$ 为图中 $\alpha$，$r$ 为图中 $\beta$），由几何关系得 $\sin i=\frac{d}{\sqrt{d^{2}+h^{2}}}$，$\sin r=\frac{d}{\sqrt{d^{2}+H^{2}}}$，因 $d\ll h$，$d\ll H$，故 $\frac{H^{2}}{h^{2}}=\left(\frac{4}{3}\right)^{2}$，得 $H=\frac{4}{3}h$。

\onepicture[w=4cm]{figs/20_an.jpg}
}

\item
%% number: 21
%% paperName: 1990年上海高考第 21 题 4 分
%% typeId: 3
%% type: 填空题
%% chapter: 运动和力
%% point: 连接体
%% method: 无
%% score: 4
%% degree: 650
%% duplicateId: 0
%% body:
如图所示的三个物体质量分别为 $m_{1}$、$m_{2}$ 和 $m_{3}$，带有滑轮的物体放在光滑水平面上，滑轮和所有接触面的摩擦以及绳子的质量均不计。为使三个物体无相对运动，水平推力 $F$ 等于 \tkanswer[3cm]{$\frac{m_{2}(m_{1}+m_{2}+m_{3})g}{m_{1}}$}。
\onepicture[align=right, w=5cm]{\includesvg{figs/21}}
%% answer: $\frac{m_{2}(m_{1}+m_{2}+m_{3})g}{m_{1}}$
\jdanswer{
$\frac{m_{2}(m_{1}+m_{2}+m_{3})g}{m_{1}}$
}
%% memo:
\memoanswer{
隔离 $m_{2}$ 分析，有 $T=m_{2}g$；隔离 $m_{1}$ 分析，有 $T=m_{1}a$；对三个物体整体分析，有 $F=(m_{1}+m_{2}+m_{3})a$，联立解得 $F=\frac{(m_{1}+m_{2}+m_{3})m_{2}g}{m_{1}}$。
}

\item
%% number: 22
%% paperName: 1990年上海高考第 22 题 3 分
%% typeId: 4
%% type: 实验题
%% chapter: 气体性质
%% point: 验证玻意耳定律
%% method: 无
%% score: 3
%% degree: 650
%% duplicateId: 0
%% body:
在《验证玻意耳$--$马略特定律》的实验中，若实验时的大气压强 $p_{0}=1.00\times10^{5}\UPa$，测得活塞和框架的重 $G_{0}=0.58\UN$，活塞面积 $S=2.0\Ucmq$。把一段空气柱封闭在注射器内，用弹簧秤竖直上提活塞，测得弹簧秤上的读数 $F=3.58\UN$，则空气柱的压强 $p=$ \tkanswer[2.5cm]{$8.5\times10^{4}\UPa$}。
%% answer: $8.5\times10^{4}\UPa$
\jdanswer{
$8.5\times10^{4}\UPa$
}
%% memo:
\memoanswer{
由平衡条件得 $p_{0}S+G_{0}=F+pS$，解得 $p=p_{0}+\frac{G_{0}-F}{S}$，代入数据得 $p=8.5\times10^{4}\UPa$。
}

\item
%% number: 23
%% paperName: 1990年上海高考第 23 题 6 分
%% typeId: 4
%% type: 实验题
%% chapter: 功和能
%% point: 机械能守恒定律
%% method: 无
%% score: 6
%% degree: 650
%% duplicateId: 0
%% body:
为进行《验证机械能守恒定律》的实验，有下列器材可供选择：铁架台，打点计时器以及复写纸、纸带，低压直流电源，天平，秒表，导线，电键。其中不必要的器材是 \tkanswer[3.5cm]{低压直流电源、天平、秒表}；缺少的器材是 \tkanswer[4.2cm]{低压交流电源、重物（重锤）、刻度尺}。
%% answer: 低压直流电源、天平、秒表；低压交流电源、重物（重锤）、刻度尺
\jdanswer{
低压直流电源、天平、秒表；低压交流电源、重物（重锤）、刻度尺
}
%% memo:
\memoanswer{
利用打点计时器验证机械能守恒定律时，动能表达式为 $\frac{1}{2}mv^{2}$，重力势能的变化量为 $mgh$，因此需要重物，但不需要测量重物的质量。实验时利用交流电源给打点计时器供电，并利用刻度尺测量纸带上的点距，再利用纸带上的点距计算不同点的速度大小，因此不必要的器材是低压直流电源、天平、秒表；缺少的器材是低压交流电源、重物（重锤）、刻度尺。
}

\item
%% number: 24
%% paperName: 1990年上海高考第 24 题 6 分
%% typeId: 4
%% type: 实验题
%% chapter: 电路
%% point: 测电动势与内阻
%% method: 无
%% score: 6
%% degree: 650
%% duplicateId: 0
%% body:
用伏安法测电阻，当被测电阻的阻值不能估计时，可采用试接的办法。如图所示，让伏特表一端接在电路上的 a 点，另一端先后接到 b 点和 c 点，注意观察两个电表的示数。若安培表的示数有显著变化，则待测电阻的阻值跟 \tkanswer{伏特} 表的内阻可比拟，伏特表应接在 a、\tkanswer{c} 两点。若伏特表的示数有显著变化，则待测电阻的阻值跟 \tkanswer{安培} 表的内阻可比拟，伏特表应接在 a、\tkanswer{b} 两点。
\onepicture[align=right, w=5cm]{\includesvg{figs/24}}
%% answer: 伏特，c，安培，b
\jdanswer{
伏特，c，安培，b
}
%% memo:
\memoanswer{
若电流表的示数有显著变化，说明电压表上的电流不能忽略，则待测电阻的阻值跟电压表的内阻可比拟，应该用电流表内接法，故电压表应接在 a、c 两点；若电压表的示数有显著变化，说明电流表上的电压不能忽略，则待测电阻的阻值跟电流表的内阻可比拟，应该用电流表外接法，故电压表应接在 a、b 两点。
}

\item
%% number: 25
%% paperName: 1990年上海高考第 25 题 4 分
%% typeId: 4
%% type: 实验题
%% chapter: 光学
%% point: 透镜
%% method: 无
%% score: 4
%% degree: 650
%% duplicateId: 0
%% body:
某同学在测量凸透镜焦距的实验中，将光屏放在离物 $60\Ucm$ 处，发现无论将透镜放于物与光屏之间的什么位置，都不能在屏上形成物的像。由此可以判定此透镜的焦距\xzanswer{B}
\fourchoices[answer=B]
{一定大于 $30\Ucm$}
{一定大于 $15\Ucm$}
{一定小于 $30\Ucm$}
{一定小于 $15\Ucm$}
%% answer: B
\jdanswer{
B
}
%% memo:
\memoanswer{
由题意知 $u+v=60$，由透镜成像公式得 $\frac{1}{u}+\frac{1}{v}=\frac{1}{f}$，联立得 $u^{2}-60u+60f=0$，因为只能成虚像，所以方程无解，即判别式 $\Delta=b^{2}-4ac=60^{2}-4\times60f<0$，得 $f>15\Ucm$，故 B 正确。
}

\item
%% number: 26
%% paperName: 1990年上海高考第 26 题 4 分
%% typeId: 4
%% type: 实验题
%% chapter: 电路
%% point: 电容
%% method: 无
%% score: 4
%% degree: 650
%% duplicateId: 0
%% body:
如图所示，让平行板电容器带电后，静电计的指针偏转一定角度，若不改变 A、B 两极板带的电量而减少两极板间距离，同时在两极板间插入电介质，那么静电计指针的偏转角度\xzanswer{A}
\onepicture[align=right, w=4cm]{\includesvg{figs/26}}
\fourchoices[answer=A]
{一定减少}
{一定增大}
{一定不变}
{可能不变}
%% answer: A
\jdanswer{
A
}
%% memo:
\memoanswer{
由平行板电容器的决定式知 $C=\frac{\varepsilon S}{4\pi kd}$，当 $d$ 减小，或 $\varepsilon$ 增大时 $C$ 增大；又 $U=\frac{Q}{C}$，$Q$ 不变，所以 $U$ 减小，即静电计指针的偏转角度减小，故 A 正确。
}

\item
%% number: 27
%% paperName: 1990年上海高考第 27 题 10 分
%% typeId: 6
%% type: 计算题
%% chapter: 气体性质
%% point: 气体连接体
%% method: 无
%% score: 10
%% degree: 650
%% duplicateId: 0
%% body:
如图所示，一密闭容器内贮有一定质量的气体，不导热的光滑活塞将容器分隔成左右两部分。开始时，两部分气体的体积、温度和压强都相同，均为 $V_{0}$、$T_{0}$ 和 $p_{0}$。将两边气体加热到某一温度，而右边仍保持原来温度，平衡时，测得右边气体的压强为 $p$，求左边气体的温度 $T$。
\onepicture[align=right, w=4cm]{\includesvg{figs/27}}
%% answer: $T=\left(\frac{2p}{p_{0}}-1\right)T_{0}$
\jdanswer{
$T=\left(\frac{2p}{p_{0}}-1\right)T_{0}$
}
%% memo:
\memoanswer{
设重新平衡时左、右两边的体积为 $V_{1}$、$V_{2}$；左边压强为 $p_{1}$，由理想气体状态方程，对左边气体：$\frac{p_{1}V_{1}}{T}=\frac{p_{0}V_{0}}{T_{0}}$。

对右边气体 $T_{0}$ 不变，由玻意耳定律得：$pV_{2}=p_{0}V_{0}$。

由压强平衡得 $p_{1}=p$，由几何关系得 $V_{1}+V_{2}=2V_{0}$。

联立解得 $T=\left(\frac{2p}{p_{0}}-1\right)T_{0}$。
}

\item
%% number: 28
%% paperName: 1990年上海高考第 28 题 13 分
%% typeId: 6
%% type: 计算题
%% chapter: 电场
%% point: 带电粒子的偏转
%% method: 无
%% score: 13
%% degree: 450
%% duplicateId: 0
%% body:
一质量为 $m$、带电量为 $+q$ 的小球从距地面高为 $h$ 处以一定的初速度水平抛出。在距抛出点水平距离为 $l$ 处，有一根管口比小球直径略大的竖直细管，管的上口距地面 $\frac{h}{2}$。为使小球能无碰撞地通过管子，可在管子上方的整个区域里加一个场强方向水平向左的匀强电场，如图所示，求：
\begin{enumerate}
	\item
	小球的初速度 $v_{0}$；
	\item
	电场强度 $E$ 的大小；
	\item
	小球落地时的动能 $E_{k}$。
\end{enumerate}
\onepicture[align=right, w=4.5cm]{\includesvg{figs/28}}
%% answer: （1）$v_{0}=2l\sqrt{\frac{g}{h}}$；（2）$E=\frac{2mgl}{qh}$；（3）$E_{k}=mgh$
\jdanswer{
\begin{enumerate}
	\item
	$v_{0}=2l\sqrt{\frac{g}{h}}$

	\item
	$E=\frac{2mgl}{qh}$

	\item
	$E_{k}=mgh$

\end{enumerate}
}
%% memo:
\memoanswer{
（1）小球运动至管上口的时间由竖直方向的运动决定，由运动学公式得 $\frac{h}{2}=\frac{1}{2}gt^{2}$。

为使带电小球无碰撞地通过细管，到达管上口的速度方向应竖直向下，即水平速度降为零，设水平方向加速度为 $a$，有 $v_{0}=at$，$v_{0}^{2}=2al$，联立解得 $v_{0}=2l\sqrt{\frac{g}{h}}$。

（2）由（1）知 $a=\frac{v_{0}}{t}=v_{0}\sqrt{\frac{g}{h}}$，由牛顿第二定律得 $qE=ma$，联立解得 $E=\frac{2mgl}{qh}$。

（3）为求小球落地时的动能 $E_{k}$，由动能定理：$E_{k}-\frac{1}{2}mv_{0}^{2}=mgh-qEl$，以 $v_{0}$、$E$ 的值代入，得 $E_{k}=mgh$。
}

\item
%% number: 29
%% paperName: 1990年上海高考第 29 题 15 分
%% typeId: 6
%% type: 计算题
%% chapter: 电磁感应
%% point: 电磁感应能量分析
%% method: 无
%% score: 15
%% degree: 450
%% duplicateId: 0
%% body:
如图 \ref{1990上海29a}，边长为 $l$ 和 $L$ 的矩形线圈 aa$'$、bb$'$ 互相垂直，彼此绝缘，可绕中心轴 O$_{1}$O$_{2}$ 转动。将两线圈的始端并在一起接到滑环 C，末端并在一起接到滑环 D，C、D 绝缘。电阻 $R=2r$，通过电刷跟 C、D 连接。线圈处于磁铁和圆柱形铁心之间的磁场中，磁场边缘对中心的张角为 $45^{\circ}$，如图 \ref{1990上海29b} 所示[图 \ref{1990上海29b} 中的圆表示圆柱形铁心，它使磁铁和铁心间的磁场沿半径方向，如图中箭头所示]。不论线圈转到磁场中的什么位置，磁场的方向总是沿着线圈平面，磁场在长为 $l$ 的边所在处的磁感应强度大小恒为 $B$。设线圈 aa$'$ 和 bb$'$ 的电阻都是 $r$，两个线圈以角速度 $\omega$ 逆时针匀速转动。
\begin{enumerate}
	\item
	求线圈 aa$'$ 转到图 \ref{1990上海29b} 所示位置时，感生电动势的大小。
	\item
	求转动过程中电阻 $R$ 上的电压最大值。
	\item
	从线圈 aa$'$ 进入磁场开始计时，正确作出 $0-T$（$T$ 是线圈转动周期）时间内通过 $R$ 的电流强度 $i_{R}$ 随时间 $t$ 变化的图象[画在图 \ref{1990上海29c} 上]。
	\item
	求外力驱动两线圈转动一圈所做的功。
\end{enumerate}
\threepicture[labelA=1990上海29a, labelB=1990上海29b, labelC=1990上海29c, wA=4cm, wB=4.3cm, wC=4.3cm, align=right]
{\includesvg{figs/29a}}
{\includesvg{figs/29b}}
{\includesvg{figs/29c}}
%% answer: （1）$E=BlL\omega$；（2）$U_{R}=\frac{2}{5}BlL\omega$；（3）如图；（4）$W=\frac{3\pi B^{2}L^{2}l^{2}\omega}{5r}$
\jdanswer{
\begin{enumerate}
	\item
	$E=BlL\omega$

	\item
	$U_{R}=\frac{2}{5}BlL\omega$

	\item
	如图

	\item
	$W=\frac{3\pi B^{2}L^{2}l^{2}\omega}{5r}$

\end{enumerate}
}
%% memo:
\memoanswer{
（1）由磁场特点知，线圈不论转到磁场中哪个位置，切割速度始终与磁场方向垂直，切割产生的电动势为 $E=2Blv=2Bl\cdot\frac{\omega}{2}=BlL\omega$。

（2）线圈 aa$'$ 或 bb$'$ 在磁场中时，等效电路如图示，电源内电阻为 $r$，外电路总电阻为 $R_{\text{外}}=\frac{Rr}{R+r}=\frac{2}{3}r$，两端电压为 $U_{R}=IR_{\text{外}}=\frac{E}{r+\frac{2}{3}r}\cdot\frac{2}{3}r=\frac{2}{5}E=\frac{2}{5}BlL\omega$。3 分。未画等效电路，计算正确同样给分。

\onepicture[w=5cm]{\includesvg{figs/29_an2}}

（3）线圈 aa$'$ 或 bb$'$ 在磁场中时，通过 $R$ 的电流大小相等。$i_{R\text{最大}}=\frac{U_{R}}{R}=\frac{\frac{2}{5}BlL\omega}{2r}=\frac{BlL\omega}{5r}$。从线圈 aa$'$ 进入磁场开始计时，每隔 $\frac{T}{8}$（线圈转动 $45^{\circ}$），电流发生一次改变，$i_{R}$ 随时间 $t$ 的变化图象如图所示。

\onepicture[w=7cm]{\includesvg{figs/29_an}}

（4）$i_{\text{总}}=\frac{E}{r+\frac{2}{3}r}=\frac{3E}{5r}$，外力驱动两线圈转动一周所做的功等于同一时间内电流做的功，$W=Pt=4Ei_{\text{总}}\cdot\frac{T}{8}=\frac{3E^{2}}{5r}\cdot\frac{T}{2}=\frac{3(BlL\omega)^{2}}{5r}\cdot\frac{\pi}{\omega}=\frac{3\pi B^{2}L^{2}l^{2}\omega}{5r}$。
}

\end{enumerate}

\end{document}
